\subsection{The remaining cyclic-four comparison must reverse the pairing}

\begin{proposition}[B-RD40: no direct nine-congruence in the cyclic-four branch]
\label{prop:B-RD40}
\status{Proved. Complete root and independent proof, original-source, and exact-certificate audits passed.}
Let $K$ be as in Proposition~\ref{prop:B-RD38}, and let $E,E'/K$
have common irreducible cyclic-four residual representation. If there
is a symplectic isomorphism $E[9]\simeq E'[9]$, their requested
prime-to-three special-inertia characters agree.

Consequently any opposite-character level-$27$ comparison in the
$C_4$ row of (B-RD39.3) has pairing multiplier congruent to $-1$
modulo three. Every such comparison can be made antisymplectic by a
scalar change; none can be made symplectic. This is an orientation
restriction, not an exclusion of the entire $C_4$ row.
\end{proposition}
\begin{proof}
Put $L=K(E[3])=K(E'[3])$ and let $F$ be its quadratic subfield.
Use the given symplectic identification to choose symplectic bases of
both three-torsion groups with the same actual $C_4$ Galois matrix.
As in Proposition~\ref{prop:B-RD38}, choose the bases so that its
projective action on $X(3)$ is
\[
 h(t)=\frac{t+2}{t-1},\qquad \sigma(t)=h(t),
                                                        \tag{B-RD40.1}
\]
where $\sigma\in\operatorname{Gal}(F/K)$ is nontrivial.
Write $t_E,t'$ for the two parameters, $z=t_E-1$, $z'=t'-1$,
and choose $\pi^2=3$ in $K$. Thus
\[
 u=z/\pi,\quad u'=z'/\pi,\qquad
 N_{F/K}(u)=N_{F/K}(u')=1,\qquad
 \overline u,\overline{u'}\in\{1,-1\}.
                                                        \tag{B-RD40.2}
\]

\paragraph{The actual marked differential cocycle.}
We determine the lift of (B-RD40.1), rather than using only its
projective action. Let $\zeta$ be a primitive cube root and
$c=\zeta-\zeta^2$, so $c^2=-3$. On the Hesse pencil
\[
 F_t(x,y,z)=x^3+y^3+z^3-3txyz
\]
use the matrix
\[
 M=\begin{pmatrix}1&1&1\\1&\zeta&\zeta^2\\
                         1&\zeta^2&\zeta\end{pmatrix}.
\]
Direct polynomial calculation gives
\[
 \det M=-3c,\qquad
 F_t(MX)=3(1-t)F_{h(t)}(X).
                                                        \tag{B-RD40.3}
\]
For the Poincare residue differential
$\eta_t=\operatorname{Res}(\Omega/F_t)$, where
$\Omega=x\,dy\wedge dz-y\,dx\wedge dz+z\,dx\wedge dy$, this gives
\[
                  M^*\eta_t=\kappa(t)\eta_{h(t)},
        \qquad    \kappa(t)=\frac c{t-1},
        \qquad    \kappa(t)\kappa(h(t))=-1.
                                                        \tag{B-RD40.4}
\]
The two Hesse translation generators are the cyclic coordinate
permutation and $\operatorname{diag}(1,\zeta,\zeta^2)$.
Neither generator has a fixed point on a smooth member of the
pencil: the diagonal generator fixes only coordinate vertices, and
the eigenvectors of the cyclic permutation lie on the cubic only
when $t=1$. Thus both are translations; they generate all nine
three-torsion translations. Conjugation by $M$ interchanges them
with one inverted, so its marked action is an order-four symplectic matrix, whose square is
$-I$. Moving the image flex back to the chosen origin is a
translation and does not change the invariant differential.
The possible central sign is fixed by the actual common marked
matrix on $E[3]$ and $E'[3]$, hence is the same on both sides.
The use of these marked Hesse translations agrees with the full-level
moduli convention in \citet[Section~2, pp.~4--5]{Fisher2015Nine}.

We also check that passing to the Weierstrass family used in
Fisher's cube equations introduces no parameter-dependent factor.
On the plane cubic set
\[
              X=z,\qquad Y=x-y,\qquad Z=x+y+tz.
\]
Its equation becomes
\[
 3ZY^2=4(t^3-1)X^3-9t^2X^2Z+6tXZ^2-Z^3.
                                                        \tag{B-RD40.5}
\]
Put $a=4(t^3-1)/3$ and take affine coordinates $aX/Z,aY/Z$,
followed by the usual $x$-translation to remove the square term.
The resulting short Weierstrass coefficients are
\[
 A_H=-\frac{t(t^3+8)}3,\qquad
 B_H=\frac{2(t^6-20t^3-8)}{27}.
\]
The displayed family $E_t:y^2+3txy+(t^3-1)y=x^3$ instead has
\[
 A_E=-\frac{3t(t^3+8)}{16},\qquad
 B_E=-\frac{t^6-20t^3-8}{32}.
\]
They are related by the constant Weierstrass scaling $2/c$.
The determinant of the preceding plane coordinate change is two;
tracking the residue differential through (B-RD40.5) gives
\[
                         \eta_t=\frac c3\,\omega_t,
                                                        \tag{B-RD40.6}
\]
where $\omega_t$ is the invariant differential of $E_t$ under this
identification. Thus (B-RD40.4) is also the exact marked differential
factor for this Weierstrass family. This computation supplies the
normalization behind the constant twist mentioned in
\citet[Lemma~4.2(i), p.~8]{Fisher2015Nine}.

\paragraph{Different tame characters force opposite unit residues.}
Take marked isomorphisms $E_L\simeq E_{t_E,L}$ and
$E'_L\simeq E_{t',L}$, and let their pullbacks of the universal
invariant differentials be $\alpha\omega_E$ and
$\alpha'\omega_{E'}$, respectively. For a generator
$\tau\in\operatorname{Gal}(L/K)$, equation (B-RD40.4) and the
common actual marked matrix show that
\[
       \frac{\tau(\alpha'/\alpha)}{\alpha'/\alpha}
                =\frac{z'}z
          \quad\hbox{or its inverse},
                                                        \tag{B-RD40.7}
\]
according to the simultaneous choice of generator and direction of
the universal isomorphism. Either choice has the same residue for
the argument below. A common central minus sign cancels in the ratio.

Let $\delta,\delta'$ be discriminant valuations over $K$, and
put $n=v_L(\alpha)$, $n'=v_L(\alpha')$. The invariant discriminant
section $\Delta\omega^{12}$ and the valuation
$v_L(\Delta(E_t))=60$ from Proposition~\ref{prop:B-RD38} give
\[
                 4\delta=60+12n,\qquad
                 4\delta'=60+12n'.
                                                        \tag{B-RD40.8}
\]
If the tame characters differ, their orders are one and two, so
$\delta'-\delta\equiv2\pmod4$. Hence
$n'-n\equiv2\pmod4$. In a tame cyclic extension of degree four,
the residue of $\tau(\beta)/\beta$ for an element of valuation
$m$ is $\xi^m$, where $\xi$ is a primitive fourth root.
Apply this to $\beta=\alpha'/\alpha$. Equation (B-RD40.7) forces
\[
                           \overline{u'}=-\overline u.
                                                        \tag{B-RD40.9}
\]

\paragraph{The direct nine-congruence contradicts a norm cube.}
Over $F$, twist both curves by the same quadratic character of $L/F$.
Their three-torsion becomes rational, they become the marked curves
$E_{t_E},E_{t'}$, and the symplectic nine-congruence persists.
Fisher's direct equation (10) therefore makes
\[
       \frac{t'-\theta}{t_E-\theta}\in(F^\times)^3
          \qquad(\theta=1,\zeta,\zeta^2).
                                                        \tag{B-RD40.10}
\]
At $\theta=1$ this says $u'/u$ is a cube. Its negative is a cube
congruent to one by (B-RD40.9). Since $v_F(3)=4$, the expansion
$(1+w)^3-1=3w+3w^2+w^3$ shows that any such cube has deviation
valuation at least three. Thus $v_F(u'+u)\ge3$.

Set
\[
 T=z+3/z=\pi(u+u^{-1}),\qquad
 T'=z'+3/z'=\pi(u'+u'^{-1})\in K.
\]
Both $T$ and $T'$ have $K$-valuation one. Moreover
\[
       T'+T=\pi(u'+u)\bigl(1+(u'u)^{-1}\bigr).
\]
The second factor has positive $F$-valuation by (B-RD40.9), and
$v_F(\pi)=2$, so $v_K(T'+T)\ge3$.
For $\lambda=\zeta-1$, the exact relation
$\lambda^2=-3\zeta$ gives
\[
 N_{F/K}(t_E-\zeta)
       =(z-\lambda)(3/z-\lambda)=-\lambda(T+3).
\]
Taking the norm of the cube in (B-RD40.10) at $\theta=\zeta$
makes
\[
                 q=\frac{T'+3}{T+3}\in(K^\times)^3.
\]
But
\[
              -q-1=-\frac{T'+T+6}{T+3}
\]
has $K$-valuation exactly one: its numerator has valuation two
and its denominator valuation one. This contradicts
Lemma~\ref{lem:B-RD37-units}, since $-q$ is a cube congruent to one.
The assumed different tame characters are impossible.

Finally, a unit of $\mathbf Z/27$ congruent to one modulo three is
a square, by Hensel lifting for $X^2-a$. A torsion isomorphism with
such a multiplier could be rescaled to be symplectic, contradicting
the result just proved on nine-torsion. Every remaining multiplier
is congruent to $-1$, and its negative is a square, giving the stated
antisymplectic normalization.
\end{proof}

\paragraph{Exact certificate and scope.}
The program \path{code/verify_c4_orientation_gate.py} checks all
Fourier, cubic, differential-scaling and short-Weierstrass identities,
the exact norm identity, the valuation inequalities, and every
unit multiplier modulo twenty-seven. Its retained output records
the author runtime $0.1905790330056334$ seconds. An initial
certificate used structural equality for two differently factored
SymPy expressions; that failed version and traceback are retained in
\path{failed_candidates/B40_certificate_structural_equality_v1.py} and
its adjacent output. The corrected test proves their exact difference
is zero. No antisymplectic level-$27$ pair is constructed or
excluded by this proposition alone; it is excluded in
Theorem~\ref{thm:B-RD42}. Independent author-output-preserving reruns by the
root and second reviewers passed in $0.1678152059976128$ and
$0.1795537689977209$ seconds, respectively.
